% =============================================================================
% Semana 11 — Dashboards que Forçam Decisão + KPIs + Pitch de 3 minutos
% Aula de 3 horas (19h–22h) | 13/05/2026
% =============================================================================

\chapter{Semana 11 (13/05/2026) — Dashboards que Forçam Decisão e o Pitch de 3 Minutos}

% ---------------------------------------------------------------------------
\section{Objetivo da semana}
% ---------------------------------------------------------------------------

Pessoal, vocês já sabem criar visuais claros e interativos. Hoje a pergunta é:
\textbf{o dashboard de vocês realmente força uma decisão?} Ou ele apenas
\emph{exibe informação} sem deixar claro o que o usuário deve fazer?

Um dashboard pode ser tecnicamente impecável e, ainda assim, não provocar
nenhuma ação. O problema geralmente é de \textbf{estrutura narrativa} e
de \textbf{clareza de KPIs}. E no final da aula, cada um vai apresentar
a Big Idea do seu T6 em exatamente 3 minutos.

\begin{BoardBox}
\textbf{Roteiro da aula (19h–22h)}
\begin{enumerate}
  \item 19h–20h — \textbf{TEORIA:} Estrutura do dashboard decisório
  \item 20h–20h40 — \textbf{TEORIA:} KPIs, targets e indicadores visuais de status
  \item 20h40–21h — \textbf{TEORIA:} O Pitch de 3 minutos — estrutura e prática
  \item 21h–21h40 — \textbf{PRÁTICA:} Revisão em par do T6 (peer review estruturado)
  \item 21h40–22h — \textbf{Pitches rápidos:} Big Idea de cada aluno (1 min cada)
\end{enumerate}
\end{BoardBox}

% =============================================================================
\section{Parte I — Estrutura do Dashboard Decisório (19h–20h)}
% =============================================================================

\begin{BoardBox}
\textbf{As três zonas de um dashboard decisório (para o quadro)}

\textbf{Zona 1 — Contexto:} O que está acontecendo agora?
\begin{itemize}
  \item KPIs principais com comparação ao período anterior ou à meta.
  \item Nenhuma interpretação ainda — apenas ``o estado atual''.
  \item Deve ser lida em $<$ 5 segundos (Teste do Relance para KPIs).
\end{itemize}

\textbf{Zona 2 — Diagnóstico:} Por que está acontecendo?
\begin{itemize}
  \item Drill-down dos KPIs em alerta: por produto, região, período, segmento.
  \item Permite formular hipóteses, não confirmar causas (isso é AED!).
  \item Filtros ou ações de dashboard levam o usuário aqui após ver a Zona 1.
\end{itemize}

\textbf{Zona 3 — Recomendação:} O que fazer?
\begin{itemize}
  \item Síntese das evidências em 1–2 recomendações concretas.
  \item Pode ser um painel de texto estático ou um visual de simulação
  (ex.: parâmetro de meta que mostra ``se aumentarmos a meta em X\%, o impacto seria Y'').
  \item O Call to Action deve ser claro: \emph{quem} faz \emph{o quê} até \emph{quando}.
\end{itemize}
\end{BoardBox}

\begin{SolvedBox}
\textbf{Dashboard decisório vs dashboard informativo — diferença crítica:}

\begin{tabular}{|l|p{6cm}|p{5.5cm}|}
\hline
& \textbf{Informativo} & \textbf{Decisório} \\
\hline
Objetivo & Reportar o que aconteceu & Provocar uma ação específica \\
KPI sem meta & Apenas o número & Número + meta + gap \\
Título & ``Vendas do Mês'' & ``Vendas 18\% abaixo da meta — risco de fechamento'' \\
Estrutura & Todos os dados disponíveis & Apenas o que suporta a decisão \\
Call to Action & Ausente & Explícito na última tela \\
\hline
\end{tabular}
\end{SolvedBox}

\begin{BoardBox}
\textbf{Anti-padrões de dashboard (erros clássicos a evitar)}
\begin{itemize}
  \item \textbf{Tudo em verde:} toda KPI sempre verde independente do contexto.
  Cria conforto falso e impede decisões urgentes.
  \item \textbf{Métricas sem definição:} ``Receita'' pode ser bruta, líquida,
  recorrente mensal, anual. Sem definição explícita, cada pessoa lê diferente.
  \item \textbf{Granularidade excessiva:} 50 KPIs na tela inicial. Onde o usuário
  deve olhar primeiro?
  \item \textbf{Falta de comparação:} um número isolado não tem significado.
  R\$100.000 de vendas é ótimo ou péssimo? Comparado ao quê?
  \item \textbf{Período sem contexto:} ``Jan'' — Jan de qual ano? É comparação YoY?
  MTD? YTD?
\end{itemize}
\end{BoardBox}

% =============================================================================
\section{Parte II — KPIs, Targets e Indicadores de Status (20h–20h40)}
% =============================================================================

\begin{SolvedBox}
\textbf{Anatomia de um bom KPI em dashboard:}

Um KPI sozinho não decide nada. Um KPI bem construído tem 5 componentes:
\begin{enumerate}
  \item \textbf{O valor atual:} o número real do período.
  \item \textbf{A meta (target):} o que deveria ser — benchmark, histórico ou meta corporativa.
  \item \textbf{O gap:} distância entre o atual e a meta (em valor absoluto e \%).
  \item \textbf{A tendência:} está melhorando ou piorando? (seta ou sparkline)
  \item \textbf{O período:} MTD (mês até hoje), YTD (ano até hoje), vs ano anterior.
\end{enumerate}

\textbf{Em Tableau:} KPI card com \texttt{IF [Atual] $\geq$ [Meta] THEN "\texttt{◆}" ELSE "\texttt{▼}" END}
como indicador de status (verde/vermelho via cor condicional).
\end{SolvedBox}

\begin{BoardBox}
\textbf{Tipos de visualização para métricas e targets (para o quadro)}

\begin{tabular}{|l|p{6.5cm}|p{4cm}|}
\hline
\textbf{Visual} & \textbf{Quando usar} & \textbf{No Tableau} \\
\hline
KPI card & Meta de valor único (receita, NPS, churn) & Text viz + forma \\
Bullet chart & Comparar atual vs meta vs bom/satisfatório & Gantt reproposto \\
Gauge & Progresso percentual para audiência não técnica & Gráfico de rosca \\
Sparkline & Tendência rápida sem detalhes de eixo & Linha em mini-tela \\
Linha de referência & Meta no gráfico de linhas ou barras & Linha de referência calculada \\
\hline
\end{tabular}

\textbf{Bullet chart} (Stephen Few) é o mais eficiente para comparação meta vs atual:
uma barra representa o valor real; outra, mais fina e mais escura, a meta.
Mais informação por centímetro quadrado do que qualquer gauge.
\end{BoardBox}

\begin{SolvedBox}
\textbf{Consistência de definições — a causa invisível de dashboards inúteis:}

Antes de construir qualquer KPI, documente:
\begin{itemize}
  \item \textbf{Nome exato:} ``Receita Líquida MRR'' — não ``Receita''.
  \item \textbf{Cálculo:} fórmula ou regra de negócio completa.
  \item \textbf{Período:} qual janela temporal e como atualiza (diário, semanal).
  \item \textbf{Fonte:} de qual tabela/sistema o dado vem.
  \item \textbf{Dono:} quem valida se o número está errado.
\end{itemize}

Sem isso, dois analistas apresentando ``Receita'' na mesma reunião podem
apresentar números diferentes — e nenhum estará errado; estarão calculando
coisas diferentes.
\end{SolvedBox}

% =============================================================================
\section{Parte III — O Pitch de 3 Minutos (20h40–21h)}
% =============================================================================

\begin{BoardBox}
\textbf{Estrutura do pitch de 3 minutos (para o quadro)}

3 minutos $\approx$ 450 palavras $\approx$ 4 frases por bloco.

\begin{enumerate}
  \item \textbf{Situação (30 seg):} contexto factual que a audiência reconhece.
  ``Nossa taxa de churn no Q1 foi de 22\% — acima dos 15\% históricos.''
  \item \textbf{Complicação (45 seg):} o problema ou oportunidade que a situação revela.
  ``Ao segmentar por canal, descobrimos que 80\% do churn vem de clientes
  adquiridos via campanha de redes sociais no mesmo período.''
  \item \textbf{Resolução (60 seg):} a evidência e a recomendação.
  ``O custo de retenção via programa de onboarding é R\$15/cliente — vs
  R\$180 para adquirir um novo. Recomendo redirecionar 30\% do orçamento
  de aquisição para retenção dos primeiros 90 dias.''
  \item \textbf{Próximo passo (45 seg):} pedido concreto à audiência.
  ``Preciso da aprovação de vocês para testar esse programa em 500 clientes
  no próximo mês e medir o impacto antes de escalar.''
\end{enumerate}
\end{BoardBox}

\begin{SolvedBox}
\textbf{Checklist do pitch eficaz:}
\begin{itemize}
  \item[$\square$] Começo com dado, não com contexto genérico
  \item[$\square$] Big Idea está na primeira frase da complicação
  \item[$\square$] A recomendação tem número concreto (não ``melhorar'' — sim ``aumentar em 15\%'')
  \item[$\square$] O próximo passo tem: verbo + quem + prazo
  \item[$\square$] Sem jargão técnico não essencial para o público
  \item[$\square$] Cronometrado: paro quando bate 3 minutos, não quando termino
\end{itemize}
\end{SolvedBox}

% =============================================================================
\section{Parte IV — Peer Review do T6 (21h–21h40)}
% =============================================================================

\begin{BoardBox}
\textbf{Roteiro de peer review estruturado (em par — 20 min cada):}

O revisor usa este checklist enquanto o colega apresenta o dashboard:

\begin{tabular}{|p{5cm}|l|p{5.5cm}|}
\hline
\textbf{Critério} & \textbf{Sim/Não} & \textbf{Sugestão específica} \\
\hline
Big Idea clara na primeira tela? & & \\
Todos os títulos são mensagem? & & \\
Passa no Teste do Relance? & & \\
KPIs têm meta + gap? & & \\
Interatividade funciona? & & \\
Estrutura: contexto→diag→reco? & & \\
Call to action explícito? & & \\
\hline
\end{tabular}

\textbf{Após o checklist,} o revisor dá 1 feedback de melhoria prioritária
e 1 ponto forte genuíno. O aluno apresentado anota e tem tempo para ajustar
antes da Semana 12.
\end{BoardBox}

% =============================================================================
\section{Parte V — Pitches Rápidos (21h40–22h)}
% =============================================================================

\begin{SolvedBox}
Cada aluno apresenta a Big Idea do seu T6 em \textbf{1 minuto}:
\begin{itemize}
  \item 15 seg: situação + problema
  \item 30 seg: evidência central
  \item 15 seg: recomendação + call to action
\end{itemize}

A turma anota: ``Ficou clara a mensagem? O que faltou?''
Feedback coletivo de 1 minuto após cada pitch.
\end{SolvedBox}

% ---------------------------------------------------------------------------
\section{Materiais Preparados}
% ---------------------------------------------------------------------------
\begin{itemize}
  \item Slides: anatomia do KPI card; bullet chart (como construir no Tableau)
  \item Formulário de peer review (impresso, 1 por par)
  \item Cronômetro visível para os pitches
  \item Exemplos de dashboards decisórios reais (Tableau Public gallery)
\end{itemize}

% ---------------------------------------------------------------------------
\section{Próximo passo}
% ---------------------------------------------------------------------------

Na \textbf{Semana 12} é a velika final: Apresentações do Trabalho 6
(pitch completo de 3 min + demo do dashboard) + revisão para a Prova II.
Cheguem com o dashboard no Tableau Public publicado e o link testado.

Perguntem aos alunos: \textit{``Imaginem que o gerente da empresa que vocês
estão analisando vai ver o dashboard pela primeira vez sem nenhuma explicação
de vocês. Qual é a primeira ação que ele tomaria com base no que está na tela?
Se vocês não souberem responder, o dashboard ainda não está pronto.''}
